\section{Implicaciones del sistema y gobernanza futura}

\subsection{Panel de implicaciones del sistema}
Hubert llamó al diseño de AlphaProof «arquitectura interpretable de tres capas»: 1) Formalizer/Proof Generator/Verifier conforman el deterministic pipeline; 2) Trace/log + dashboard constituyen la observability layer; 3) Governance review + newsletter son la human-in-the-loop layer. La experiencia de AlphaProof nos dice que un agent complejo debe separar con claridad en capas el pipeline, la metric y la gobernanza para poder desplegarse de forma segura en un ámbito sin matices de gris como math.

\begin{table}[H]
\centering
\begin{tabular}{p{3cm}p{4cm}p{5cm}}
\toprule
Nivel & Punto de atención & Método \\
\midrule
Pipeline & correctness + modularity & Formalizer/Proof Generator/Verifier + Runner conforman steeped pipeline \\
Observabilidad & metrics & Dashboard + chunk gating + trace log \\
Governance & trust & human review + newsletter + failure curriculum \\
\bottomrule
\end{tabular}
\caption{Las tres capas de implicaciones del sistema de AlphaProof}
\end{table}

\begin{knowledgebox}{Esencia de las implicaciones del sistema}
Concebir el system design como una arquitectura de tres capas — pipeline + observabilidad + governance — permite garantizar que la tarea de razonamiento matemático más sensible mantenga su interpretabilidad y su control incluso en escenarios desconocidos.
\end{knowledgebox}

\subsection{Dirección de la gobernanza futura}
Hubert propuso la idea de «differentiated governance»: fijar un drift threshold distinto para cada search track y registrar las governance decision en el newsletter/backlog. Por ejemplo, ante una vulnerabilidad de beam search, el governance reviewer aplica de manera directa un patch a la policy correspondiente, mientras que en el heuristic track es el RL agent quien reintenta de forma autónoma. Él enfatiza que la gobernanza no es un judge, sino una feedback policy.

\begin{warningbox}{La gobernanza no debe reducirse a prohibir}
Cuando se activa una governance alert, no se trata de desactivar directamente al agent, sino de escribir de vuelta el human comment en el failure curriculum, aumentar el reward shaping, y después reanudar la search; de lo contrario, el system terminará convertido en un fragile checkbox.
\end{warningbox}

\subsection{Resumen de la sección}
Las implicaciones del sistema se resumen en una arquitectura de tres capas, mientras que la gobernanza futura, mediante differential threshold + feedback loop, permite que el agent matemático siga expandiendo sus límites; estas reflexiones también pueden retroalimentar otros sistemas RL que requieren un alto grado de interpretabilidad.
\newpage

